\documentclass[10pt]{article}

% ---------- Document Identity (update these only) ----------
\newcommand{\DocStage}{}
\newcommand{\DocTitle}{Your Road to Tier One SOF}
\newcommand{\ProgramName}{182-Week Civilian-to-Selection Readiness Pipeline}
\newcommand{\ProgramSubtitle}{}

% ---------- Page ----------
\usepackage[legalpaper,left=0.5in,right=0.5in,top=0.6in,bottom=0.45in]{geometry}

% ---------- Typography ----------
\usepackage[T1]{fontenc}
\usepackage[utf8]{inputenc}
\usepackage{libertinus}
\usepackage{microtype}
\usepackage{parskip}
\usepackage{ragged2e}
\usepackage{textcomp}
\AtBeginDocument{\color{black}}
\AtBeginDocument{\pagecolor{white}}
\AtBeginDocument{\justifying}

% ---------- Landscape pages ----------
\usepackage{pdflscape}

% ---------- Color ----------
\usepackage[table]{xcolor}
\definecolor{StageBlue}{HTML}{1F4E79}
\definecolor{StageBlueDark}{HTML}{163A5A}
\definecolor{LightGray}{HTML}{F2F4F7}
\definecolor{MidGray}{HTML}{D9DEE5}
\definecolor{RuleGray}{HTML}{C7CED8}
\definecolor{HeaderInk}{HTML}{000000}
\definecolor{Ink}{HTML}{000000}

% ---------- Utility ----------
\usepackage{enumitem}
\sloppy
\setlength{\emergencystretch}{2em}

% ---------- Boxes / UI ----------
\usepackage[most]{tcolorbox}

\tcbset{
  charterTitle/.style={
    enhanced,
    colback=StageBlue!4,
    colframe=StageBlue,
    boxrule=1.25pt,
    arc=3pt,
    left=16pt,right=16pt,top=14pt,bottom=14pt,
    borderline north={3.2pt}{0pt}{StageBlue},
    borderline west={4pt}{0pt}{StageBlue},
    borderline south={1.25pt}{0pt}{StageBlue},
    drop shadow={black!25!white},
  },
  charterRule/.style={
    enhanced,
    colback=LightGray,
    colframe=RuleGray,
    boxrule=0.85pt,
    arc=3pt,
    left=12pt,right=12pt,top=10pt,bottom=10pt,
    borderline west={3pt}{0pt}{StageBlue},
  },
  charterBadge/.style={
    enhanced,
    colback=StageBlue,
    coltext=white,
    colframe=StageBlueDark,
    boxrule=0.6pt,
    arc=2pt,
    left=6pt,right=6pt,top=3pt,bottom=3pt,
  }
}
\newtcbox{\Badge}{nobeforeafter,charterBadge}

% ---------- Lists ----------
\setlist[itemize]{leftmargin=1.5em, itemsep=0.25em, topsep=0.25em}
\setlist[enumerate]{leftmargin=1.8em, itemsep=0.25em, topsep=0.25em}

% ---------- TOC ----------
\usepackage{tocloft}
\setcounter{tocdepth}{3}
\setlength{\cftbeforesecskip}{3pt}
\setlength{\cftbeforesubsecskip}{2pt}
\setlength{\cftbeforesubsubsecskip}{0pt}
\renewcommand{\cftsecfont}{\bfseries}
\renewcommand{\cftsecpagefont}{\bfseries}
\renewcommand{\cftsubsecfont}{\normalfont}
\renewcommand{\cftsubsecpagefont}{\normalfont}
\renewcommand{\cftsubsubsecfont}{\normalfont}
\renewcommand{\cftsubsubsecpagefont}{\normalfont}
\setlength{\cftsecindent}{0pt}
\setlength{\cftsubsecindent}{1.25em}
\setlength{\cftsubsubsecindent}{2.8em}
\setlength{\cftsecnumwidth}{2.1em}
\setlength{\cftsubsecnumwidth}{2.8em}
\setlength{\cftsubsubsecnumwidth}{3.6em}

% Remove the built-in TOC heading so only the custom heading is shown
\makeatletter
\renewcommand{\tableofcontents}{\@starttoc{toc}}
\makeatother

% ---------- Header/Footer: page number only ----------
\usepackage{fancyhdr}
\pagestyle{fancy}
\fancyhf{}
\renewcommand{\headrulewidth}{0pt}
\renewcommand{\footrulewidth}{0pt}
\fancyfoot[C]{\thepage}

% ---------- Section formatting ----------
\usepackage{titlesec}

\titleformat{\section}
  {\Large\bfseries\color{Ink}}
  {\thesection}{0.6em}{}
  [\vspace{0.25em}\textcolor{StageBlue}{\rule{\linewidth}{0.8pt}}]

\titleformat{\subsection}
  {\large\bfseries\color{Ink}}
  {\thesubsection}{0.6em}{}

\titleformat{\subsubsection}
  {\normalsize\bfseries\color{Ink}}
  {\thesubsubsection}{0.6em}{}

\titlespacing*{\section}{0pt}{0.95em}{0.35em}
\titlespacing*{\subsection}{0pt}{0.65em}{0.25em}
\titlespacing*{\subsubsection}{0pt}{0.55em}{0.2em}

% ---------- Table engine ----------
\usepackage{tabularray}
\UseTblrLibrary{booktabs}

% ---------- tabularray: GLOBAL table treatment ----------
\SetTblrInner{
  cells = {fg=black, bg=white, valign=t},
}

% ---------- Header helpers ----------
\newcommand{\Hdr}[1]{\shortstack[c]{\strut #1\strut}}

% ---------- Lines ----------
\newcommand{\TitleLine}{\textcolor{StageBlue}{\rule{0.98\linewidth}{0.95pt}}}
\newcommand{\SubTitleLine}{\textcolor{RuleGray}{\rule{0.98\linewidth}{0.65pt}}}

% ---------- Continued on next page ----------
\newcommand{\ContinuedOnNextPageInline}{%
  \par
  \vfill
  \noindent\textcolor{RuleGray}{\rule{\linewidth}{1pt}}\vspace{-2pt}
  \noindent\hfill\textit{Continued on next page}\par\vspace{-10pt}
  \noindent\textcolor{RuleGray}{\rule{\linewidth}{1pt}}
  \clearpage
}

% ---------- PDF hyperlinks ----------
\usepackage[hidelinks]{hyperref}
\hypersetup{
  colorlinks=false,
  pdfborder={0 0 0},
  linktoc=all,
  pdftitle={\DocTitle},
  pdfauthor={\ProgramName}
}

% =========================================================
\begin{document}

\subsection*{Phase 13 Card Header}
\phantomsection
\addcontentsline{toc}{subsection}{Phase 13 Card Header}

\begin{itemize}
\item Phase \textbf{ID}: Phase 13
\item Phase Name: Full-Spectrum Readiness Consolidation
\item Duration weeks: 12
\item Version: v1.0
\item Stage Alignment: Stage 5
\item Governing Status: Draft — Non-Deployable
\end{itemize}

\subsection*{1. Phase Mission}
\phantomsection

Consolidate final, low-variance verification across selection-relevant domains using admissible evidence produced only inside protected windows, closing replication gaps and demonstrating stable durability and recovery posture without unsafe stacking.

\subsection*{2. Phase Purpose}
\phantomsection

\begin{itemize}
\item Purpose in sequence: provide the final consolidation and proof phase where the program emphasizes repeatability, configuration freeze, and audit-ready evidence chains over further expansion of training stress.
\item Setup for next phase: this is the terminal consolidation phase; its output is a final readiness evidence packet (tests + logs) that supports Stage 9 packaging and deployable implementation posture.
\end{itemize}

\subsection*{3. Entry Criteria}
\phantomsection

\begin{itemize}
\item Entry requires admissible prior-phase exit posture per Stage 1.F (\textbf{VALID} sessions and \textbf{VALID} tests) and Phase 12 completion consistent with Stage 3 reslot discipline.
\item Default posture if evidence is incomplete or invalid: \textbf{HOLD} and reslot missing/invalid gate tests into the next protected window per Stage 3 adjustment doctrine.
\end{itemize}

\subsection*{4. Operating Constraints}
\phantomsection

\begin{itemize}
\item 6 sessions per week.
\item Required session elements in correct order:
  \begin{itemize}
  \item Safety self-check first
  \item Warm-up
  \item Mobility
  \item Main work
  \item Cardio
  \item Cooldown
  \item Recovery notes and any required post-session notes per Stage 1.F
  \end{itemize}
\item Cardio minimum standard: minimum 30 minutes continuous per session.
\item 10,000 steps per day global requirement.
\item Validity doctrine: admissible \textbf{VALID} evidence only for gates; \textbf{INVALID} provides no gate credit and triggers reslot/\textbf{HOLD} posture per Stage 3.
\end{itemize}

\newpage
\subsection*{5. Primary Adaptations and Physiological Focus}
\phantomsection

\begin{itemize}
\item Reduce variance of outputs: demonstrate repeatable performance under stable environment/tool conditions and stable recovery posture.
\item Close replication gaps for gate-critical metrics: ensure tier classification evidence meets replication requirements under Stage 2 posture.
\item Maintain durability constraints: preserve gait stability, foot integrity, and symptom-free tolerance across verification weeks.
\item Preserve specificity: prioritize selection-relevant tests (\textbf{RUN}, \textbf{RUCK}, \textbf{CAL}, \textbf{SWIM}) without adding new instruments or unauthorized proxies.
\item Maintain \textbf{BODYCOMP} directionality stability without measurement drift (method consistency and complete logging).
\end{itemize}

\subsection*{6. Primary Modalities}
\phantomsection

\begin{itemize}
\item Emphasized domains: \textbf{RUN}, \textbf{RUCK}, \textbf{CAL}, \textbf{SWIM}, \textbf{DURABILITY}, \textbf{RECOVERY}, \textbf{BODYCOMP}.
\item Supporting domains: \textbf{STR} (submax proxy), \textbf{AER} (non-run conditioning, non-interchangeable with \textbf{RUN}).
\item De-emphasized/prohibited:
  \begin{itemize}
  \item No new tests, no new domains, no new thresholds.
  \item No unsupervised underwater/breath-hold; underwater is governance-gated and conditional only under Stage 2 constraints.
  \item No gate-grade tests outside protected windows.
  \end{itemize}
\end{itemize}

\clearpage
\begin{landscape}

\subsection*{7. Weekly Structure}
\phantomsection

\begingroup
\scriptsize
\setlength{\tabcolsep}{2pt}
\renewcommand{\arraystretch}{1.12}

\begin{longtblr}{
  width=\linewidth,
  rowhead=1,
  colspec = {
    Q[l,wd=0.90in]
    X[l,1.55]
    Q[l,wd=1.20in]
    X[l,2.85]
    X[l,2.10]
  },
  hlines,
  vlines,
  cells = {hsep=6pt, vsep=6pt, halign=l, valign=t},
  cell{1-Z}{1-5} = {fg=black},
  row{odd} = {bg=white},
  row{even} = {bg=LightGray},
  row{1} = {bg=StageBlue!15, fg=black, font=\bfseries, halign=l, valign=t},
}
\Hdr{Session \#} &
\Hdr{Primary Focus} &
\Hdr{Main Work Domains} &
\Hdr{Cardio Modality/Intent} &
\Hdr{Notes/Risk Controls} \\

1 &
Locomotion readiness and durability control &
\textbf{DURABILITY}; \textbf{RUN} &
Modality: walk/jog as phase-appropriate; Minutes: ≥30; RPE plus Talk Test: RPE 3–5, full to short sentences; Continuity: continuous planned &
Protect gait; no novelty near test windows; any mechanics-altering pain triggers downshift and documentation. \\

2 &
Calisthenics maintenance and movement quality &
\textbf{CAL}; \textbf{DURABILITY} &
Modality: lowest-risk steady modality; Minutes: ≥30; RPE plus Talk Test: RPE 3–4, full sentences; Continuity: continuous planned &
\textbf{CAL} artifact posture applies on test weeks; stop before form degradation; no failure chasing. \\

3 &
Ruck interface stability and durability support &
\textbf{RUCK}; \textbf{DURABILITY} &
Modality: low-impact steady (walk-based); Minutes: ≥30; RPE plus Talk Test: RPE 3–5, full to short sentences; Continuity: continuous planned &
Footwear/pack interface is controlled; avoid stacking with maximal \textbf{RUN} within spacing intent; reslot if foot flags rise. \\

4 &
Swim readiness and recovery protection &
\textbf{SWIM}; \textbf{RECOVERY} &
Modality: low-risk steady modality; Minutes: ≥30; RPE plus Talk Test: RPE 3–4, full sentences; Continuity: continuous planned &
Pool access and unit verification required for gate-grade swim tests; reslot if pool conditions compromise validity. \\

5 &
Strength proxy support or recovery-biased maintenance &
\textbf{STR}; \textbf{RECOVERY} &
Modality: lowest-risk steady modality; Minutes: ≥30; RPE plus Talk Test: RPE 2–4, full sentences; Continuity: continuous planned &
\textbf{STR} remains submax proxy only; do not let \textbf{STR} contaminate \textbf{RUN}/\textbf{RUCK} verification windows; prioritize recovery stability. \\

6 &
Consolidation and pre-window stabilization &
\textbf{RECOVERY}; \textbf{DURABILITY} &
Modality: low-risk steady modality; Minutes: ≥30; RPE plus Talk Test: RPE 2–4, full sentences; Continuity: continuous planned &
Treat as stabilization day before protected windows; no novelty; reinforce logging completeness. \\

\end{longtblr}

\endgroup

\end{landscape}
\clearpage

\subsection*{8. Progression Rules}
\phantomsection

\begin{itemize}
\item Phase 13 progression is verification-first, not expansion-first:
  \begin{itemize}
  \item do not increase exposure complexity or intensity to chase new performance;
  \item instead, reduce variance and improve repeatability under controlled conditions.
  \end{itemize}
\item Anti-stacking safeguard: if multiple major tests are planned in a protected week, prioritize the decisive gate items and reslot the rest rather than compressing.
\item Regression triggers: any red-flag symptom, mechanics-altering pain, foot breakdown trend, or recovery collapse triggers downshift to Minimum Viable Sessions and reslot of tests.
\item Substitution posture: preserve intent, reduce risk, and maintain admissibility; substitutions that break comparability are treated as non-gate-grade and require retest under valid conditions.
\end{itemize}

\subsection*{9. Deload and Test Placement}
\phantomsection

\begin{itemize}
\item Phase 13 taper-aware posture per Stage 3 integrated calendars:
  \begin{itemize}
  \item Week 1 baseline.
  \item Week 5 Deload+Test.
  \item Week 10 Deload+Test.
  \item Week 11 Transition.
  \item Week 12 Deload+Test end gate.
  \end{itemize}
\item Gate-grade tests occur only during Deload+Test weeks; Transition week is stabilization and retest buffer only under validity-preserving conditions.
\end{itemize}

\begin{landscape}

\newpage
\subsection*{10. Test Battery}
\phantomsection

\begingroup
\scriptsize
\setlength{\tabcolsep}{2pt}
\renewcommand{\arraystretch}{1.12}

\begin{longtblr}{
  width=\linewidth,
  rowhead=1,
  colspec = {
    Q[l,wd=1.35in]
    X[l,2.35]
    Q[l,wd=0.80in]
    X[l,3.10]
    Q[l,wd=1.40in]
  },
  hlines,
  vlines,
  cells = {hsep=6pt, vsep=6pt, halign=l, valign=t},
  cell{1-Z}{1-5} = {fg=black},
  row{odd} = {bg=white},
  row{even} = {bg=LightGray},
  row{1} = {bg=StageBlue!15, fg=black, font=\bfseries, halign=l, valign=t},
}
\Hdr{Test Timing} &
\Hdr{Test \textbf{ID}/Name} &
\Hdr{Domain} &
\Hdr{Validity Conditions} &
\Hdr{Pass Threshold Stage 2 Tier} \\

Week 5 Deload+Test &
\textbf{C7} / \textbf{MET-2B-013} Plank — Forearm &
\textbf{CAL} &
\textbf{VALID} only; auditable start/stop; form termination enforced; required fields logged &
\textbf{SELECTION-READY} \\

Week 5 Deload+Test &
\textbf{C13} / \textbf{MET-2B-023} Swim 500 m Time — No Fins &
\textbf{SWIM} &
\textbf{VALID} only; pool unit verified (meters); no fins; allowable strokes; rest governance honored; required fields logged &
\textbf{SELECTION-READY} \\

Week 10 Deload+Test &
\textbf{C12} / \textbf{MET-2B-021} 12-Mile Ruck Time Trial — 55 lb &
\textbf{RUCK} &
\textbf{VALID} only; dry load verified; water excluded from official load and logged separately; route measured or treadmill identity/settings logged per protocol; no stop-rule violations &
\textbf{SELECTION-READY} \\

Week 10 Deload+Test &
\textbf{C16} / \textbf{MET-2B-036} Gait-Altering Pain Flag &
\textbf{DURABILITY} &
Must be No for gate posture; logged contemporaneously; no inference &
\textbf{SELECTION-READY} \\

Week 12 Deload+Test End Gate &
\textbf{C11} / \textbf{MET-2B-017} 5-Mile Run Time Trial &
\textbf{RUN} &
\textbf{VALID} only; route measured or treadmill identity/settings logged per protocol; no stop-rule violations; environment and tool stability logged &
\textbf{SELECTION-READY} \\

Week 12 Deload+Test End Gate &
\textbf{C4} / \textbf{MET-2B-010} Push-Ups — 2-Minute Timed, Strict Standard &
\textbf{CAL} &
\textbf{VALID} only; \textbf{CAL} artifact required for gate-validity; strict standard; required fields logged &
\textbf{SELECTION-READY} \\

Week 12 Deload+Test End Gate &
\textbf{C5} / \textbf{MET-2B-011} Sit-Ups — 2-Minute Timed, Standard &
\textbf{CAL} &
\textbf{VALID} only; \textbf{CAL} artifact required for gate-validity; required fields logged &
\textbf{SELECTION-READY} \\

Week 12 Deload+Test End Gate &
\textbf{C6} / \textbf{MET-2B-012} Pull-Ups — Strict Dead-Hang &
\textbf{CAL} &
\textbf{VALID} only; \textbf{CAL} artifact required for gate-validity; pronated strict standard; required fields logged &
\textbf{SELECTION-READY} \\

Week 12 Deload+Test End Gate &
\textbf{C13} / \textbf{MET-2B-022} Swim 500 yd Time — No Fins &
\textbf{SWIM} &
\textbf{VALID} only; pool unit verified (yards); no fins; allowable strokes; rest governance honored &
\textbf{SELECTION-READY} \\

Week 12 Deload+Test End Gate &
\textbf{C14} / \textbf{MET-2B-025} Underwater Exposure Admissibility Flag &
\textbf{SWIM} &
Governance-gated: requires certified lifeguard + dedicated safety spotter; if governance absent, not scheduled and inadmissible &
\textbf{SELECTION-READY} \\

Week 12 Deload+Test End Gate &
\textbf{C14} / \textbf{MET-2B-026} Underwater 25 m Time — One Breath &
\textbf{SWIM} &
Governance-gated: only if \textbf{MET-2B-025} is Yes and all supervision prerequisites are met; otherwise not scheduled and inadmissible &
\textbf{SELECTION-READY} \\

\end{longtblr}

\endgroup

\end{landscape}
\clearpage

\subsection*{11. Exit Standards}
\phantomsection

\begin{itemize}
\item Phase 13 exit posture is \textbf{SELECTION-READY} theme, supported by admissible \textbf{VALID} evidence and replication posture per Stage 2.
\item Gate outcomes are applied per Stage 1.F: \textbf{ADVANCE}/\textbf{HOLD}/\textbf{REGRESS}/\textbf{MEDICAL} \textbf{HOLD}.
\item Any \textbf{INVALID} test provides no gate credit and requires reslot per Stage 3; no advancement claims are made on invalid evidence.
\item Underwater is optional and governance-gated; it cannot be used for gate credit if governance prerequisites are not met and is not replaced by a proxy.
\end{itemize}

\subsection*{12. Risk Controls and Stop Rules}
\phantomsection

\begin{itemize}
\item Stage 0.5 and Stage 1.F govern stop posture and escalation; Stage 3.E governs reslot/extension posture.
\item Phase 13 risk controls (phase-specific emphasis):
  \begin{itemize}
  \item protect \textbf{RUN} and \textbf{RUCK} from fatigue contamination by spacing intent and staged batteries; reslot rather than stack.
  \item enforce strict environment/tool stability (route \textbf{ID}, pool verification, treadmill/machine identity) to preserve admissibility.
  \item enforce foot integrity and gait stability as primary disqualifiers (\textbf{MET-2B-036}).
  \item underwater is never attempted without governance prerequisites; if prerequisites are absent, it is not scheduled and is inadmissible.
  \end{itemize}
\end{itemize}

\clearpage
\begin{landscape}

\subsection*{13. Required Capabilities and Dependencies}
\phantomsection

\begingroup
\scriptsize
\setlength{\tabcolsep}{2pt}
\renewcommand{\arraystretch}{1.12}

\begin{longtblr}{
  width=\linewidth,
  rowhead=1,
  colspec = {
    X[l,1.55]
    X[l,2.60]
    Q[l,wd=1.05in]
    Q[l,wd=1.00in]
    X[l,2.80]
  },
  hlines,
  vlines,
  cells = {hsep=6pt, vsep=6pt, halign=l, valign=t},
  cell{1-Z}{1-5} = {fg=black},
  row{odd} = {bg=white},
  row{even} = {bg=LightGray},
  row{1} = {bg=StageBlue!15, fg=black, font=\bfseries, halign=l, valign=t},
}
\Hdr{Dependency \textbf{ID}} &
\Hdr{Required Capability Category and Tier} &
\Hdr{Due-By week-in-phase} &
\Hdr{Owner} &
\Hdr{Fail Action} \\

\textbf{DEP-1C-006} &
7.18 Testing Timing Measurement and Course-Setup Tools — Tier: required &
Week 1 &
Equipment Lead &
\textbf{HOLD} gate-grade \textbf{RUN}/\textbf{RUCK} tests until route/course setup is verifiable; reslot to next protected window; treat any unverified attempt as \textbf{INVALID} for gate use. \\

\textbf{DEP-1C-007} &
7.11 Footwear Systems and Foot Care — Tier: required &
Week 1 &
Trainee &
\textbf{HOLD} \textbf{RUN}/\textbf{RUCK} gate attempts if foot breakdown trend exists; substitute to lowest-risk Cardio modality for sessions; reslot tests until stable. \\

\textbf{DEP-1C-013} &
7.07 Load-Bearing Rucking and Carry Systems — Tier: required &
Week 5 &
Equipment Lead &
\textbf{HOLD} ruck gate attempts; do not improvise load verification; reslot to next protected window; invalid attempts provide no gate credit. \\

\textbf{DEP-1C-011} &
7.17 Aquatic Training and Water Safety Equipment — Tier: required &
Week 5 &
Coach Lead &
If pool access or unit verification is not stable, reslot \textbf{SWIM} tests; do not substitute non-swim evidence for \textbf{SWIM} gate claims. \\

\textbf{DEP-1C-017} &
7.12 Recovery Mobility and Tissue-Care Implements — Tier: required &
Week 1 &
Trainee &
If durability drift rises, downshift to Minimum Viable Sessions and reslot tests; do not stack make-up testing. \\

\textbf{DEP-1C-015} &
7.09 Health Monitoring Diagnostics and Biometrics — Tier: supporting &
Week 1 &
Coach Lead &
If monitoring is unavailable, proceed with conservative defaults; do not infer missing data; missing data is recorded as \textbf{MISSING}/\textbf{UNKNOWN} and may constrain gate confidence. \\

\textbf{DEP-1C-016} &
7.10 Logistics Maintenance Storage and Sustainment Equipment — Tier: supporting &
Week 1 &
Equipment Lead &
If sustainment posture fails, correct before test windows to prevent evidence loss; preserve Evidence Ref capture for all tests. \\

\end{longtblr}

\endgroup

\end{landscape}
\clearpage

\subsection*{14. Validity Requirements}
\phantomsection

\begin{itemize}
\item Evidence admissibility is governed by Stage 1.F and Stage 2:
  \begin{itemize}
  \item tests are \textbf{VALID}/\textbf{INVALID} only;
  \item missing required fields or required environment/tool identity defaults to \textbf{INVALID} for gate use.
  \end{itemize}
\item \textbf{CAL} tests require compliant artifact evidence for \textbf{VALID} gate use.
\item Replication posture applies for tier classification and gate confidence; gate reliance requires replicated \textbf{VALID} outcomes per Stage 2 doctrine.
\item \textbf{INVALID} tests provide no gate credit and require reslot per Stage 3; do not compress by stacking.
\end{itemize}

\subsection*{15. Change Control Notes}
\phantomsection

\begin{itemize}
\item Allowed without patch: reslotting within protected windows and re-ordering staged batteries within a Deload+Test week when it preserves Stage 3 anti-stacking intent and does not change the test set.
\item Requires patch: changing gate test set, adding protected windows, altering phase duration, altering domain emphasis materially, or introducing new dependencies as hard requirements.
\item Missed/invalid tests are handled per Stage 3.E reslot discipline; do not stack make-up testing to meet deadlines.
\end{itemize}

\subsection*{16. Compliance Checklist}
\phantomsection

\begin{itemize}
\item Phase \textbf{ID}, name, duration, and governing status stated.
\item Operating constraints include required session elements, Cardio minimum, and steps requirement in body text.
\item Weekly Structure table includes six sessions and Cardio cell structure (Modality; Minutes; RPE plus Talk Test; Continuity continuous planned).
\item Deload/test windows identified and aligned to Stage 3 phase calendar posture.
\item Test Battery lists only Stage 2 tests/metrics (C\# and \textbf{MET-2B}-###) and uses tier labels only.
\item Dependencies are category-level using 7.01–7.20; 7.02 is not required.
\item Governing Status set to Draft — Non-Deployable because dependency IDs are placeholders (“use Stage 1 dependency \textbf{ID}”) in gate-relevant dependency fields.
\end{itemize}

\end{document}